\chapter{Lesson 7: Telephone and Appointments}

Phone calls are hard in a foreign language because there are no gestures to help. This
lesson collects fixed telephone expressions, how to read phone numbers, the progressive
\ko{\til 고 있다} (``What are you doing right now?''), and the two reason connectors
\ko{\til 아서/어서} and \ko{\til (으)니까}, which we needed for excuses (``I'm late
because \ldots'').

\section{Vocabulary and telephone expressions}

\begin{table}[H]
\centering
\begin{tabular}{lll}
\toprule
Korean & RR & English \\
\midrule
\ko{전화} & jeonhwa & telephone, phone call \\
\ko{전화번호} & jeonhwabeonho & phone number \\
\ko{휴대폰 / 핸드폰} & hyudaepon / haendeupon & cell phone \\
\ko{전화하다 / 전화를 걸다} & jeonhwahada / jeonhwareul geolda & to call \\
\ko{전화를 받다} & jeonhwareul batda & to answer the phone \\
\ko{문자(를) 보내다} & munja(reul) bonaeda & to send a text \\
\ko{약속을 잡다} & yaksogeul japda & to set up an appointment \\
\ko{약속을 바꾸다 / 취소하다} & yaksogeul bakkuda / chwisohada & to change / cancel plans \\
\ko{늦다} & neutda & to be late \\
\ko{일찍} & iljjik & early \\
\ko{오전 / 오후} & ojeon / ohu & a.m. / p.m. \\
\ko{반} & ban & half (past the hour) \\
\ko{잠깐만} & jamkkanman & just a moment \\
\bottomrule
\end{tabular}
\end{table}

\begin{table}[H]
\centering
\begin{tabular}{lll}
\toprule
Expression & RR & Meaning \\
\midrule
\ko{여보세요?} & Yeoboseyo? & Hello? (on the phone only) \\
\ko{거기 김민수 씨 댁이지요?} & Geogi Gim Minsu ssi daegijiyo? & Is this Kim Minsu's home? \\
\ko{민수 씨 좀 바꿔 주세요.} & Minsu ssi jom bakkwo juseyo. & Could I speak to Minsu? \\
\ko{전데요.} & Jeondeyo. & Speaking. (lit.\ it's me) \\
\ko{잠깐만 기다리세요.} & Jamkkanman gidariseyo. & Please hold on a moment. \\
\ko{지금 안 계신데요.} & Jigeum an gyesindeyo. & He/she isn't here right now. \\
\ko{메시지 남겨 드릴까요?} & Mesiji namgyeo deurilkkayo? & Shall I take a message? \\
\ko{전화 잘못 거셨어요.} & Jeonhwa jalmot geosyeosseoyo. & You have the wrong number. \\
\ko{다시 전화할게요.} & Dasi jeonhwahalgeyo. & I'll call back. \\
\ko{들어가세요.} & Deureogaseyo. & Goodbye (ending a call, lit.\ go in). \\
\bottomrule
\end{tabular}
\end{table}

\begin{remark}
The \ko{\til 지요?} (\ko{\til 죠?}) ending asks for confirmation, like a tag question:
\ko{댁이지요?} (this is \ldots's home, right?). The \ko{\til (으)ㄴ데요/는데요} ending in
\ko{전데요} and \ko{안 계신데요} softens a statement and leaves room for the listener to
respond; we learned these as set phrases.
\end{remark}

\subsection{Reading phone numbers}

Phone numbers use Sino-Korean digits read one at a time. Zero is \ko{공} gong (or
\ko{영} yeong), and the hyphen is read as \ko{에} (written \ko{의}, the possessive
particle, which is pronounced [\ko{에}] here).

\begin{example}[Phone numbers]
Korean mobile numbers start with 010:
\begin{itemize}
  \item 010-4385-2179: \ko{공일공의 사삼팔오의 이일칠구}, read gongilgong-e samsampalo-e
        iilchilgu.
  \item \ko{전화번호가 몇 번이에요?} Jeonhwabeonhoga myeot beonieyo? (What is your phone
        number?)
\end{itemize}
\end{example}

\section{Grammar 1: Progressive \ko{\til 고 있다}}

Stem + \ko{\til 고 있다} describes an action in progress (``be \ldots ing''). It attaches
directly to the stem, and only \ko{있다} conjugates.

\begin{center}
\begin{tikzpicture}[x=1cm, every node/.style={font=\small}]
  \draw[-{Stealth}, thick] (0,0) -- (10,0) node[right] {time};
  \fill[orange!30] (2.5,-0.25) rectangle (7.5,0.25);
  \draw (2.5,-0.25) rectangle (7.5,0.25);
  \node at (5,0) {\ko{공부하고 있어요}};
  \fill[red!60] (5,0.45) circle (2.5pt);
  \node[above, font=\scriptsize] at (5,0.5) {now: the action has started and not yet finished};
  \fill[orange!20] (0.4,-1.35) rectangle (3.0,-0.85); \draw (0.4,-1.35) rectangle (3.0,-0.85);
  \node at (1.7,-1.1) {\ko{공부하고 있었어요}};
  \node[font=\scriptsize] at (1.7,-1.6) {was studying (at that time)};
  \draw[dashed] (1.7,-0.85) -- (1.7,0);
\end{tikzpicture}
\end{center}

\begin{table}[H]
\centering
\begin{tabular}{llll}
\toprule
Form & Korean & RR & Meaning \\
\midrule
present & \ko{먹고 있어요} & meokgo isseoyo & is eating \\
past & \ko{먹고 있었어요} & meokgo isseosseoyo & was eating \\
honorific & \ko{드시고 계세요} & deusigo gyeseyo & (grandpa) is eating \\
negative & \ko{안 먹고 있어요} & an meokgo isseoyo & isn't eating \\
\bottomrule
\end{tabular}
\end{table}

\begin{example}[What are you doing?]
On the phone:
\begin{itemize}
  \item \ko{지금 뭐 하고 있어요? --- 숙제하고 있어요.} Jigeum mwo hago isseoyo? ---
        Sukjehago isseoyo. (What are you doing now? --- I'm doing homework.)
  \item \ko{어제 전화했을 때 뭐 하고 있었어요? --- 샤워하고 있었어요.} Eoje
        jeonhwahaesseul ttae mwo hago isseosseoyo? --- Syawohago isseosseoyo. (What were
        you doing when I called yesterday? --- I was taking a shower.)
  \item \ko{아버지께서 지금 운전하고 계세요.} Abeojikkeseo jigeum unjeonhago gyeseyo. (My
        father is driving right now.) --- honorific \ko{\til 고 계시다}.
  \item \ko{요즘 한국어를 배우고 있어요.} Yojeum hangugeoreul baeugo isseoyo. (These days
        I'm learning Korean.) --- an ongoing period, not just this instant.
\end{itemize}
\end{example}

\begin{remark}
With verbs of wearing (\ko{입다}, \ko{신다}, \ko{쓰다}), \ko{\til 고 있다} can mean
either ``is putting on'' or ``is wearing'' (the resulting state): \ko{민지 씨는 빨간
모자를 쓰고 있어요} (Minji is wearing a red hat). This is how people are described in
the ``who is this?'' exercises.
\end{remark}

\section{Grammar 2: \ko{\til 아서/어서} (because; and then)}

\ko{\til 아서/어서} follows vowel harmony like \ko{\til 아요/어요} (\ko{하다} $\to$
\ko{해서}). It has two meanings.

\subsection{Reason or cause}

\begin{example}[Reasons]
The first clause gives the cause:
\begin{itemize}
  \item \ko{늦잠을 자서 수업에 늦었어요.} Neutjameul jaseo sueobe neujeosseoyo. (I slept
        in, so I was late for class.)
  \item \ko{배가 아파서 학교에 못 갔어요.} Baega apaseo hakgyoe mot gasseoyo. (I
        couldn't go to school because I had a stomachache.)
  \item \ko{만나서 반가워요.} Mannaseo bangawoyo. (Nice to meet you --- lit.\ because I
        met you, I'm glad.)
  \item \ko{오늘은 일요일이라서 은행이 문을 안 열어요.} Oneureun iryoiriraseo eunhaengi
        muneul an yeoreoyo. (Today is Sunday, so the bank isn't open.) --- noun +
        \ko{(이)라서}.
\end{itemize}
\end{example}

\begin{important}[Two restrictions on reason \ko{\til 아서/어서}]
\begin{enumerate}
  \item \textbf{No tense} on the first clause: \ko{늦잠을 자서} not \ko{*잤어서}.
        The tense of the whole sentence comes from the last verb.
  \item \textbf{No commands or suggestions} in the second clause: \ko{*비가 와서 우산을
        가져가세요} is ungrammatical. Use \ko{\til (으)니까} instead (below).
\end{enumerate}
\end{important}

\subsection{Sequence (and then, there)}

With motion and other closely linked actions, \ko{\til 아서/어서} means ``do A, and
\emph{in that state / at that place} do B''. Compare with \ko{\til 고}, where the two
actions are independent.
\begin{itemize}
  \item \ko{도서관에 가서 공부했어요.} Doseogwane gaseo gongbuhaesseoyo. (I went to the
        library and studied there.)
  \item \ko{친구를 만나서 같이 저녁을 먹었어요.} Chingureul mannaseo gachi jeonyeogeul
        meogeosseoyo. (I met my friend and we had dinner together.)
  \item \ko{밥을 먹고 도서관에 갔어요.} Babeul meokgo doseogwane gasseoyo. (I ate and
        then went to the library.) --- independent actions: \ko{\til 고}.
\end{itemize}

\section{Grammar 3: \ko{\til (으)니까} (since, because)}

\ko{\til (으)니까} also gives a reason, but it is more subjective (the speaker's own
justification) and, crucially, it \emph{can} be followed by a command or suggestion. It
can also take the past tense.

\begin{itemize}
  \item Vowel stem: \ko{니까}. \ko{바쁘다} $\to$ \ko{바쁘니까}.
  \item Consonant stem: \ko{으니까}. \ko{많다} $\to$ \ko{많으니까}.
  \item \ko{ㄹ} stem: drop \ko{ㄹ}. \ko{멀다} $\to$ \ko{머니까}.
  \item Past: \ko{왔으니까}, \ko{먹었으니까}.
\end{itemize}

\begin{example}[\ko{\til (으)니까} with commands and suggestions]
These cannot use \ko{\til 아서/어서}:
\begin{itemize}
  \item \ko{비가 오니까 우산을 가져가세요.} Biga onikka usaneul gajyeogaseyo. (It's
        raining, so take an umbrella.)
  \item \ko{시간이 없으니까 택시를 탑시다.} Sigani eopseunikka taeksireul tapsida. (We
        don't have time, so let's take a taxi.)
  \item \ko{학교가 머니까 일찍 출발할까요?} Hakgyoga meonikka iljjik chulbalhalkkayo?
        (School is far, so shall we leave early?)
  \item \ko{어제 많이 먹었으니까 오늘은 조금만 먹을 거예요.} Eoje mani meogeosseunikka
        oneureun jogeumman meogeul geoyeyo. (I ate a lot yesterday, so today I'll eat
        only a little.)
\end{itemize}
\end{example}

\begin{table}[H]
\centering
\begin{tabular}{llll}
\toprule
 & \ko{\til 고} & \ko{\til 아서/어서} & \ko{\til (으)니까} \\
\midrule
meaning & and; and then & because; and then (there) & since, because \\
past on 1st clause & only when listing & no & yes \\
command/suggestion after & yes & no & yes \\
tone & neutral & objective cause; polite excuses & speaker's reasoning \\
\bottomrule
\end{tabular}
\end{table}

\begin{aside}
\textbf{Apologizing for being late.} Our instructor recommended \ko{\til 아서/어서} for
excuses to people above you, since \ko{\til (으)니까} can sound like you are justifying
yourself or blaming circumstances: \ko{늦어서 죄송합니다} (I'm sorry for being late) is
the standard phrase, never \ko{*늦으니까 죄송합니다}. Korea's culture of speed
(\ko{빨리빨리} ppallippalli, ``quickly, quickly'') means that being on time for work and
classes matters, although friends are often relaxed about 10--15 minutes. Text messages
(\ko{카톡} katok, from the KakaoTalk app) are now the default way to make and change
plans; an actual phone call is often reserved for family or urgent matters.
\end{aside}

\section{Sample phone call}

\begin{quote}
\ko{A: 여보세요? 거기 박지훈 씨 댁이지요?}\\
\ko{B: 네, 그런데요. 누구세요?}\\
\ko{A: 저는 지훈 씨 친구 에밀리예요. 지훈 씨 좀 바꿔 주세요.}\\
\ko{B: 지훈이는 지금 샤워하고 있어요. 메시지 남겨 드릴까요?}\\
\ko{A: 네, 내일 약속 시간을 바꾸고 싶어서 전화했어요. 내일 오후에 수업이 생겨서 세 시 반에
만날 수 없어요. 다섯 시가 어때요?}\\
\ko{B: 네, 그렇게 전할게요.}\\
\ko{A: 감사합니다. 들어가세요.}
\end{quote}

\noindent A: Hello? Is this Park Jihun's home? / B: Yes, it is. Who's calling? / A: I'm
Jihun's friend Emily. Could I speak to Jihun? / B: Jihun is taking a shower right now.
Shall I take a message? / A: Yes, I called because I want to change the time for
tomorrow. A class came up tomorrow afternoon, so I can't meet at 3:30. How about five?
/ B: OK, I'll pass that on. / A: Thank you. Goodbye.

\section*{Review summary}
\begin{itemize}
  \item Phone: \ko{여보세요}, \ko{\ldots 좀 바꿔 주세요}, \ko{잠깐만 기다리세요},
        \ko{전화 잘못 거셨어요}; numbers digit by digit, hyphen = \ko{에}.
  \item \ko{\til 고 있다}: action in progress; \ko{\til 고 계시다} (honorific);
        \ko{\til 고 있었어요} (was \ldots ing).
  \item \ko{\til 아서/어서}: reason (no past, no commands after) or close sequence.
  \item \ko{\til (으)니까}: reason that allows past tense and commands/suggestions.
\end{itemize}
